% !TEX root = ../../thesis.tex

In programming, a \textit{keyword} is a predefined word with a reserved meaning. Keywords can generally not be used for any other purpose, and they make up a programming language's syntax~\cite{keyword}.

In our Script Assembler, we use custom keywords in a similar way, although it may be difficult to define whether or not it is useful to refer to our parser's syntax as a \textit{programming language}. 

In this chapter, the term \textit{whitespace} is defined to refer to any amount of space or tab characters. Here, for convenience, line-breaks are chosen to specifically not fall under said definition.

\section{Robot files}

In the \gls{trp} project, the Robot Framework provides the primary scripting language that makes up the tests to be run on KeTop devices. 

A Robot file contains so-called \textit{test cases}, which are blocks commonly made up of a few lines of code~\cite{robot-framework}. Despite most parts of the syntax being irrelevant to the Script Assembler's workings, in order to get a feel for the nature of this project it may still be useful to have a look at an example of a Robot script file as is seen in Listing~\ref{lst:robot_example_script}, which is taken from Robot Framework's official home page:

\begin{lstlisting}[caption={Example code of a Robot script file~\cite{robot-framework}.}, label={lst:robot_example_script}]
*** Settings ***
Documentation     A test suite for valid login.
...
...               Keywords are imported from the resource file
Resource          keywords.resource
Default Tags      positive

*** Test Cases ***
Login User with Password
    Connect to Server
    Login User            ironman    1234567890
    Verify Valid Login    Tony Stark
    [Teardown]    Close Server Connection

Denied Login with Wrong Password
    [Tags]    negative
    Connect to Server
    Run Keyword And Expect Error    *Invalid Password    Login User    ironman    123
    Verify Unauthorised Access
    [Teardown]    Close Server Connection
\end{lstlisting}

As is made clear from the example above, some lines may be preceded by whitespace. Since the Robot Framework achieves nesting by having lines be preceded by an integer multiple of exactly four spaces certain lines of code may be syntactically erroneous or interpreted differently when that amount of whitespace is changed~\cite{robot-framework-style-guide}. This is an example of a programming convention that in this thesis will be referred to as \textit{significant indentation}. Python, for instance, works in a similar way, requiring a predefined use of leading whitespace to achieve grouping of statements~\cite{py-indentation}.

Despite Robot Framework's general ease of use, proficient support for extensibility, and versatile testing capabilities, the control over the specific execution flow is limited~\cite{robot-framework-pros-cons}. In the context of \textit{KEBA Group AG}, a single \textit{KeTop} may need to run multiple Robot script files, and a Robot script file may need to be reused for multiple KeTops. However, since it may be required for a test to be made of a single script file, it might be necessary to \textit{stitch together}, or in other words to concatenate, the contents of multiple Robot script files into one. 

Without the Script Assembler, a user may run into issues where they, for instance, would like a certain piece of code to only be included on KeTops satisfing a certain condition like being of a predefined operating system, or for certain pieces of code to be executed at a time that is either before or after the remaining parts, as the Robot Framework may not provide a useful way to achieve the required flow. 

\section{What Keywords Are}

Our project allows for custom keywords, all of which being preceded by the \verb|@| character, to be defined to act as either \textit{options} or \textit{modifiers}. In a provided file, these keywords must be defined at the start of the line, optionally preceded by indentation for nesting purposes. Lines containing such valid keywords will not be included in the concatenated output file. 

A keyword may be defined to have either no or any number of optional or required arguments. If set, these arguments may then consist of any sequence of characters without line-breaks that is prefixed by the keyword with at least one whitespace inbetween. 

The difference between options and modifiers lies in that options do not include a \textit{block body}, while modifiers do. A \textit{block body} is defined to be made up of all lines of text that, when read from top to bottom, succeed the line of the keyword, up until and not including the first line that is found that does not include at least one additional level of indentation than that of the line of the keyword, or up until the end of the file, ignoring all lines that are either empty or are solely made of whitespace. The \textit{scope} of the block body refers to the parts of the code that belong to it, as opposed to the parts of the code that do not. A level of indentation is defined to be made up of either four spaces or one tab character. 

To make up for the extra level of indentation used to define the span of the body text, the parsed body text will be stripped of that one line of indentation. Non-keyword lines, and lines that do not belong to the block body of a modifier will simply be ignored by the Script Assembler and will thus be appended into the output Robot file as-is. This means that, if the provided file does not contain anything that may be matched as a keyword, the parser should leave it as-is, aside from possible minor whitespace changes. In the case of a valid Robot file, this may only not be the case in rare exceptions, as the \verb|@| character is generally not utilized in Robot files that are to be executed on a KeTop.

\section{Examples}

\subsection{Options}

For instance, take the following code snippet, which includes the option \verb|@NO_KETOP|, which takes no arguments:

\begin{lstlisting}[caption=Example of an Option before parsing., label={lst:option_example_before_parser}]
@NO_KETOP
Install MS
\end{lstlisting}

When the line containing \verb|@NO_KETOP| is encountered, the \verb|no_ketop| flag, for which the concrete behaviour is explained in Chapter~\ref{chap:keywords-list}, is enabled for the output script. Since that specific line will not be included in the output script, assuming Listing~\ref{lst:option_example_before_parser} to be the only input to the parser, the output file may look like the following:

\begin{lstlisting}[caption={Parser result of Listing~\ref{lst:option_example_before_parser}.}, label={lst:option_example_after_parser}]
Install MS
\end{lstlisting}

\subsection{Modifiers}

Another example is using the modifier \verb|@VERSION|:

\begin{lstlisting}[caption=Example of a Modifier before parsing., label={lst:modifier_example_before_parser}]
Execute on Ketop    echo A
@VERSION BIOS > 1.0.1
    IF    ${SINGLE_MODE} is $False
        Execute on Ketop    echo B
    END
Execute on Ketop    echo C
\end{lstlisting}

Here, the line two, which is made up of \verb|@VERSION BIOS > 1.0.1|, represents a boolean condition relating to the version of the KeTop device. The block body of this condition is made up by the lines three to five. According to the specification of the \verb|@VERSION| keyword, that block body will be included only if the condition stated in its argument is satisfied.

This means that, again assuming Listing~\ref{lst:modifier_example_before_parser} to be the only input to the parser, if the condition \verb|BIOS > 1.0.1| is satisfied, the output may look like the following:

\begin{lstlisting}[caption={Parser result of Listing~\ref{lst:modifier_example_before_parser} if the condition is met.}, label={lst:modifier_example_after_parser_condition_met}]
Execute on Ketop    echo A
IF    ${SINGLE_MODE} is $False
    Execute on Ketop    echo B
END
Execute on Ketop    echo C
\end{lstlisting}

Note the removal of one level of indentation of what used to be the block body. If, instead, the case of the condition is not satisfied, the output may instead be:

\begin{lstlisting}[caption={Parser result of Listing~\ref{lst:modifier_example_before_parser} if the condition is not met.}, label={lst:modifier_example_after_parser_condition_not_met}]
Execute on Ketop    echo A
Execute on Ketop    echo C
\end{lstlisting}

\section{Recursion}

Keywords may also be defined recursively. Recursion is generally defined as being \textit{the process in which a function calls itself directly or indirectly}~\cite{recursion}. In the case of the Script Assembler, it refers to the fact that block bodies of modifiers may also include modifiers, of which the block bodies may then also include modifiers, and so forth. To show this possibility, consider the following example:

\begin{lstlisting}[caption={Example of modifiers being defined recursively.}, label={lst:recursive_modifier_example}]
@BRACKETSTART
    @OS WIN10
        @VERSION BIOS >= 1.7.0
            Upload Script To KeTop    E2Speed.ps1
\end{lstlisting}

Here, the block body of \verb|@BRACKETSTART| is made up of lines two to four, of which the block body of \verb|@OS| is made up of lines three and four, of which the block body of \verb|@VERSION| is made up of line four. Despite the order of \textit{which blocks are nested inside which} not making a difference with how all existing modifiers are specified, in general, the parser will interpret this from the outside to the innermost level of nesting, similar to how nesting is handled by most ordinary programming languages. 

This means that, if the above snippet is the only input file that contains the \verb|@BRACKETSTART| keyword, given how that keyword is specified, the part of the output script that is the so-called \textit{bracket start} will be: 

\begin{lstlisting}[caption={Bracket-start parser result for Listing~\ref{lst:recursive_modifier_example} for satisfied conditions.}, label={lst:bracket_start_output_example}]
Upload Script To KeTop    E2Speed.ps1
\end{lstlisting}

In the case of both of the conditions from the respective specifications of the \verb|@OS| and \verb|@VERSION| keywords and with each of their respective arguments being met. Otherwise, the \textit{bracket start} of the output file will be empty. 

\section{Concatenation and Output}

In general, the files get concatenated in the same order that they are given, unless keywords that change that predefined flow are used. Then, to assure the final output Robot file is valid and executable on a KeTop, the following template code will be appended before the output of the entire parsing and concatenation operations:

\begin{lstlisting}[caption={Output file template appended before inserted test cases.}, label={lst:output_template_example}]
# Template file for HMI Automation
*** Settings ***
Resource          ../../../Robot/Resources/keywords.resource
Resource          ../../../Robot/Resources/variables.resource
Suite Setup       KeTop Suite Setup
Suite Teardown    KeTop Suite Teardown
Test Setup        KeTop Test Setup
Test Teardown     KeTop Test Teardown

*** Test Cases ***
# Here the test cases are inserted
\end{lstlisting}

Finally, the resulting code will be written to a file which will then be further handled to be run either locally or on a KeTop. 